\chapter{Trait \code{IngestAdapter}}
\label{ch:ingest-adapter}

\begin{note}
Stub. El refactor P2-E introdujo un wrapper de más alto nivel
sobre \code{LogSource} + \code{LogParser}: los autores de una
fuente nueva implementan \code{IngestAdapter} en una crate hija
en lugar de editar archivos del core.
\end{note}

\section{Contrato}

\begin{lstlisting}[style=rust]
#[async_trait]
pub trait IngestAdapter: Send + Sync {
    fn adapter_id(&self) -> &str;
    fn tables(&self) -> &[&'static str];
    async fn stream_triples(
        &self,
        table: &str,
        since: Watermark,
        cancel: CancellationToken,
    ) -> Result<TripleBatchStream, SourceError>;
}

pub struct TripleBatch {
    pub triples:    Vec<ParsedTriple>,
    pub watermark:  Watermark,
    pub fetched_at: i64,
}
\end{lstlisting}

\section{Implementación genérica}

\code{LogSourceAdapter} envuelve cualquier \code{LogSource} +
\code{LogParser} en un \code{IngestAdapter}. Las 4 fuentes
existentes (Log Analytics, Defender XDR, Data Lake, AAD export)
pueden exponerse como adapters sin cambio de código fuente; el
orquestador aún las sostiene de forma concreta hoy por
compatibilidad hacia atrás, pero el seam está abierto.

\section{Propagación MVCC}

\code{TripleBatch.fetched\_at} lleva el timestamp de fetch
upstream, que es exactamente lo que el writer necesita para
llamar a \code{CompactRelation::with\_ingested\_at} en cada
arista insertada --- la tercera pata de la historia MVCC P3
(Capítulo~\ref{ch:mvcc}). Cablear esa plomería end-to-end queda
como follow-up diferido.

\begin{inthecode}
\filepath{graph\_hunter\_core/src/ingest\_adapter.rs} --- trait +
\code{LogSourceAdapter}.
\end{inthecode}
